%%%
%%% Radical "⼭"
%%%
\section*{Radical 46: ``⼭''}\addcontentsline{toc}{section}{Radical 46: ⼭}\addcontentsline{loh}{figure}{\#\#\#\# 46: ⼭}

%%%%%%%%%% 山 %%%%%%%%%%
\subsection*{山}\addcontentsline{loh}{figure}{山}

\begin{Entry}{山}{3}{⼭}[Kangxi 46]
  \begin{Phonetics}{山}{shan1}[][HSK 1]
    \definition*{s.}{Sobrenome: Shan}
    \definition[座]{s.}{colina; maciço; montanha | qualquer coisa que se assemelhe a uma montanha | arbustos nos quais os bichos"-da"-seda tecem seus casulos; referindo"-se a casulos de bicho"-da"-seda | eco; metáfora para um som muito alto}
  \end{Phonetics}
\end{Entry}

\begin{Entry}{山川}{3,3}{⼭,⼮}
  \begin{Phonetics}{山川}{shan1chuan1}[][HSK 7-9]
    \definition{s.}{montanhas e rios; paisagem}
  \end{Phonetics}
\end{Entry}

\begin{Entry}{山冈}{3,4}{⼭,⼌}
  \begin{Phonetics}{山冈}{shan1gang1}[][HSK 7-9]
    \definition[座]{s.}{colina baixa; pequeno morro}
  \end{Phonetics}
\end{Entry}

\begin{Entry}{山区}{3,4}{⼭,⼖}
  \begin{Phonetics}{山区}{shan1qu1}[][HSK 5]
    \definition[片]{s.}{área montanhosa; região montanhosa | colina; serra; montanha | distrito montanhoso}
  \end{Phonetics}
\end{Entry}

\begin{Entry}{山东}{3,5}{⼭,⼀}
  \begin{Phonetics}{山东}{shan1dong1}
    \definition*{s.}{Província de Shandong (Shantung) no nordeste da China}
  \end{Phonetics}
\end{Entry}

\begin{Entry}{山羊}{3,6}{⼭,⽺}
  \begin{Phonetics}{山羊}{shan1yang2}
    \definition{s.}{cabra | (ginástica) cavalo de salto de pequeno porte}
  \end{Phonetics}
\end{Entry}

\begin{Entry}{山西}{3,6}{⼭,⾑}
  \begin{Phonetics}{山西}{shan1xi1}
    \definition*{s.}{Província de Shanxi (Shansi) no norte da China entre Hebei e Shaanxi, abreviada como 晋 (Shansi)}
  \seealsoref{晋}{jin4}
  \end{Phonetics}
\end{Entry}

\begin{Entry}{山阴}{3,6}{⼭,⾩}
  \begin{Phonetics}{山阴}{shan1yin1}
    \definition*{s.}{Condado de Shanyin em Shuozhou, Shanxi}
    \definition{s.}{lado norte (ou sombreado) de uma montanha}
  \end{Phonetics}
\end{Entry}

\begin{Entry}{山体}{3,7}{⼭,⼈}
  \begin{Phonetics}{山体}{shan1ti3}
    \definition{s.}{forma de uma montanha}
  \end{Phonetics}
\end{Entry}

\begin{Entry}{山谷}{3,7}{⼭,⾕}
  \begin{Phonetics}{山谷}{shan1gu3}[][HSK 6]
    \definition[条,个]{s.}{vale; desfiladeiro; ravina; a área baixa e estreita entre duas montanhas geralmente tem riachos no meio}
  \end{Phonetics}
\end{Entry}

\begin{Entry}{山坡}{3,8}{⼭,⼟}
  \begin{Phonetics}{山坡}{shan1po1}[][HSK 6]
    \definition[个,座,片]{s.}{encosta; encosta da montanha; a inclinação entre o topo da montanha e o terreno plano}
  \end{Phonetics}
\end{Entry}

\begin{Entry}{山岭}{3,8}{⼭,⼭}
  \begin{Phonetics}{山岭}{shan1ling3}[][HSK 7-9]
    \definition[座,片,条]{s.}{cadeia de montanhas; cordilheira | crista; montanhas altas contínuas}
  \end{Phonetics}
\end{Entry}

\begin{Entry}{山顶}{3,8}{⼭,⾴}
  \begin{Phonetics}{山顶}{shan1ding3}[][HSK 7-9]
    \definition{s.}{topo de uma colina; cume de uma montanha; o pico de uma montanha; o topo da montanha}[他们距山顶还有100米远。===Eles ainda estavam a 100 metros do cume.]
  \end{Phonetics}
\end{Entry}

\begin{Entry}{山峰}{3,10}{⼭,⼭}
  \begin{Phonetics}{山峰}{shan1feng1}[][HSK 6]
    \definition[座,个]{s.}{pico (montanha); topo alto e pontudo da montanha}
  \end{Phonetics}
\end{Entry}

\begin{Entry}{山路}{3,13}{⼭,⾜}
  \begin{Phonetics}{山路}{shan1lu4}[][HSK 7-9]
    \definition{s.}{trilha de montanha | estrada de montanha}
  \end{Phonetics}
\end{Entry}

\begin{Entry}{山寨}{3,14}{⼭,⼧}
  \begin{Phonetics}{山寨}{shan1zhai4}[][HSK 7-9]
    \definition{s.}{fortaleza na montanha (especialmente de bandidos); aldeia fortificada na montanha; vila fortificada na colina | Humorístico: “bandido”; imitador | Figurativo: imitação (de produtos); falsificação}
  \end{Phonetics}
\end{Entry}

%%%%%%%%%% 屹 %%%%%%%%%%
\subsection*{屹}\addcontentsline{loh}{figure}{屹}

\begin{Entry}{屹}{6}{⼭}
  \begin{Phonetics}{屹}{ge1}
    \definition{s.}{verruga; pequena protuberância | pequeno monte de terra}
  \end{Phonetics}
  \begin{Phonetics}{屹}{yi4}
    \definition{adj.}{Literário: elevando como um pico de montanha; representando um pico de montanha imponente)}
  \end{Phonetics}
\end{Entry}

\begin{Entry}{屹立}{6,5}{⼭,⽴}
  \begin{Phonetics}{屹立}{yi4li4}[][HSK 7-9]
    \definition{v.}{erguer"-se imponente como um gigante; ficar ereto; erguer"-se alta e firme como o pico de uma montanha, é frequentemente usada como metáfora para algo constante e inabalável}
  \synonymref{耸立}{song3li4}
  \synonymref{挺拔}{ting3ba2}
  \synonymref{挺立}{ting3li4}
  \end{Phonetics}
\end{Entry}

%%%%%%%%%% 岁 %%%%%%%%%%
\subsection*{岁}\addcontentsline{loh}{figure}{岁}

\begin{Entry}{岁}{6}{⼭}
  \begin{Phonetics}{岁}{sui4}[][HSK 1]
    \definition{clas.}{usado para anos (de idade)}
    \definition{s.}{ano (literário) | colheita do ano (literário) | idade | tempo (literário) | ano (de idade) | ano (para as colheitas)}
  \end{Phonetics}
\end{Entry}

\begin{Entry}{岁月}{6,4}{⼭,⽉}
  \begin{Phonetics}{岁月}{sui4yue4}[][HSK 5]
    \definition[段,番]{s.}{anos; ano e mês; refere"-se a tempo em geral}
  \end{Phonetics}
\end{Entry}

\begin{Entry}{岁数}{6,13}{⼭,⽁}
  \begin{Phonetics}{岁数}{sui4shu4}[][HSK 6]
    \definition{s.}{idade; anos; a idade de uma pessoa}
  \end{Phonetics}
\end{Entry}

%%%%%%%%%% 岂 %%%%%%%%%%
\subsection*{岂}\addcontentsline{loh}{figure}{岂}

\begin{Entry}{岂}{6}{⼭}
  \begin{Phonetics}{岂}{qi3}
    \definition*{s.}{Sobrenome: Qi}
    \definition{adv.}{Litarário: expressa uma pergunta retórica, equivalente a 哪里, 怎么 e 难道}
  \seealsoref{哪里}{na3li5}
  \seealsoref{难道}{nan2dao4}
  \seealsoref{怎么}{zen3me5}
  \end{Phonetics}
\end{Entry}

\begin{Entry}{岂有此理}{6,6,6,11}{⼭,⽉,⽌,⽟}
  \begin{Phonetics}{岂有此理}{qi3you3ci3li3}[][HSK 7-9]
    \definition{expr.}{Isso é um absurdo!; Que exorbitante!; Absurdo!; Como isso pode ser assim?; Ridículo!}
    \definition{interj.}{``Isso é um absurdo!''; ``Que exorbitante!''; ``Absurdo!''; ``Como isso pode ser assim?''; ``Ridículo!''}
  \end{Phonetics}
\end{Entry}

%%%%%%%%%% 岗 %%%%%%%%%%
\subsection*{岗}\addcontentsline{loh}{figure}{岗}

\begin{Entry}{岗}{7}{⼭}
  \begin{Phonetics}{岗}{gang3}
    \definition{s.}{outeiro; monte | crista; vergão (no rosto, pele, etc.) | sentinela; posto | trabalho | batida policial}
  \end{Phonetics}
\end{Entry}

\begin{Entry}{岗位}{7,7}{⼭,⼈}
  \begin{Phonetics}{岗位}{gang3wei4}[][HSK 6]
    \definition[个,类]{s.}{posto; estação; originalmente se refere ao local guardado pelos militares e pela polícia, agora se refere a uma posição geral}
  \end{Phonetics}
\end{Entry}

%%%%%%%%%% 岛 %%%%%%%%%%
\subsection*{岛}\addcontentsline{loh}{figure}{岛}

\begin{Entry}{岛}{7}{⼭}
  \begin{Phonetics}{岛}{dao3}[][HSK 6]
    \definition[个,座]{s.}{ilha; uma massa de terra menor que um continente cercada por água}
  \end{Phonetics}
\end{Entry}

\begin{Entry}{岛屿}{7,6}{⼭,⼭}
  \begin{Phonetics}{岛屿}{dao3yu3}[][HSK 7-9]
    \definition[座,些,群]{s.}{ilha; ilhota}
  \end{Phonetics}
\end{Entry}

%%%%%%%%%% 岩 %%%%%%%%%%
\subsection*{岩}\addcontentsline{loh}{figure}{岩}

\begin{Entry}{岩}{8}{⼭}
  \begin{Phonetics}{岩}{yan2}
    \definition*{s.}{Sobrenome: Yan}
    \definition{s.}{rocha; pedra | penhasco; rochedo; picos de montanhas formados por rochas salientes | caverna (nas colinas)}
  \end{Phonetics}
\end{Entry}

\begin{Entry}{岩石}{8,5}{⼭,⽯}
  \begin{Phonetics}{岩石}{yan2shi2}[][HSK 7-9]
    \definition[块]{s.}{rocha; pedra; as rochas ígneas são agregados compostos por um ou mais minerais e constituem o principal material que forma a crosta terrestre; de acordo com sua formação, podem ser divididas em três categorias: rochas ígneas, rochas sedimentares e rochas metamórficas}
  \antonymref{石头}{shi2tou5}
  \end{Phonetics}
\end{Entry}

%%%%%%%%%% 岭 %%%%%%%%%%
\subsection*{岭}\addcontentsline{loh}{figure}{岭}

\begin{Entry}{岭}{8}{⼭}
  \begin{Phonetics}{岭}{ling3}
    \definition{s.}{cordilheira}
  \end{Phonetics}
\end{Entry}

%%%%%%%%%% 岳 %%%%%%%%%%
\subsection*{岳}\addcontentsline{loh}{figure}{岳}

\begin{Entry}{岳}{8}{⼭}
  \begin{Phonetics}{岳}{yue4}
    \definition*{s.}{As Cinco Grandes Montanhas; as cinco montanhas mais importantes na cultura tradicional chinesa simbolizam a unidade e a santidade nacional, e são frequentemente usadas em contextos religiosos, de sacrifício e culturais | Sobrenome: Yue}
    \definition{s.}{montanha alta | pais da esposa | sogro/sogra; as formas de tratamento para os pais ou tios da esposa}
  \end{Phonetics}
\end{Entry}

\begin{Entry}{岳父}{8,4}{⼭,⽗}
  \begin{Phonetics}{岳父}{yue4fu4}[][HSK 7-9]
    \definition{s.}{pai da esposa; sogro}
  \end{Phonetics}
\end{Entry}

\begin{Entry}{岳母}{8,5}{⼭,⽏}
  \begin{Phonetics}{岳母}{yue4mu3}[][HSK 7-9]
    \definition[位,名,个]{s.}{mãe da esposa, sogra}
  \end{Phonetics}
\end{Entry}

%%%%%%%%%% 岸 %%%%%%%%%%
\subsection*{岸}\addcontentsline{loh}{figure}{岸}

\begin{Entry}{岸}{8}{⼭}
  \begin{Phonetics}{岸}{an4}[][HSK 5]
    \definition{adj.}{arrogante; orgulhoso; grandioso (de maneira sombria ou condescendente)}
    \definition[条,道,段,面]{s.}{margem; costa; litoral; terreno à beira da água}
  \end{Phonetics}
\end{Entry}

\begin{Entry}{岸上}{8,3}{⼭,⼀}
  \begin{Phonetics}{岸上}{an4shang4}[][HSK 5]
    \definition{s.}{em terra; costa; margem | na margem do rio; na beira do rio}
  \end{Phonetics}
\end{Entry}

%%%%%%%%%% 峡 %%%%%%%%%%
\subsection*{峡}\addcontentsline{loh}{figure}{峡}

\begin{Entry}{峡}{9}{⼭}
  \begin{Phonetics}{峡}{xia2}
    \definition[个]{s.}{(frequentemente como parte do nome de um lugar) desfiladeiro | estreito | desfiladeiro; um lugar entre duas montanhas e um rio}
  \end{Phonetics}
\end{Entry}

\begin{Entry}{峡谷}{9,7}{⼭,⾕}
  \begin{Phonetics}{峡谷}{xia2gu3}[][HSK 7-9]
    \definition[段,条]{s.}{desfiladeiro; cânion; o rio atravessa um vale profundo e estreito, com penhascos íngremes em ambos os lados}
  \synonymref{山谷}{shan1gu3}
  \end{Phonetics}
\end{Entry}

%%%%%%%%%% 峰 %%%%%%%%%%
\subsection*{峰}\addcontentsline{loh}{figure}{峰}

\begin{Entry}{峰}{10}{⼭}
  \begin{Phonetics}{峰}{feng1}
    \definition{clas.}{usado para camelos}
    \definition{s.}{pico; cume; o pico proeminente de uma montanha | coisa parecida com um pico; coisas em forma de montanhas}
  \end{Phonetics}
\end{Entry}

\begin{Entry}{峰会}{10,6}{⼭,⼈}
  \begin{Phonetics}{峰会}{feng1hui4}[][HSK 6]
    \definition{s.}{cúpula; reunião de cúpula}
  \end{Phonetics}
\end{Entry}

\begin{Entry}{峰回路转}{10,6,13,8}{⼭,⼞,⾜,⾞}
  \begin{Phonetics}{峰回路转}{feng1hui2-lu4zhuan3}[][HSK 7-9]
    \definition{expr.}{cume em meio a elevações circundantes e estradas sinuosas;  (estrada de montanha) torcendo e virando; a estrada da montanha serpenteia em torno de cada novo pico | boa (ou nova) reviravolta nos acontecimentos; uma oportunidade surgiu inesperadamente; as coisas tomaram um novo rumo}
  \end{Phonetics}
\end{Entry}

%%%%%%%%%% 崇 %%%%%%%%%%
\subsection*{崇}\addcontentsline{loh}{figure}{崇}

\begin{Entry}{崇}{11}{⼭}
  \begin{Phonetics}{崇}{chong2}
    \definition*{s.}{Sobrenome: Chong}
    \definition{adj.}{alto; elevado; sublime}
    \definition{v.}{adorar; reverenciar; venerar; estimar | respeitar}
  \end{Phonetics}
\end{Entry}

\begin{Entry}{崇尚}{11,8}{⼭,⼩}
  \begin{Phonetics}{崇尚}{chong2shang4}[][HSK 7-9]
    \definition{v.}{sustentar; defender; valorizar}[我们崇尚公平与正义。===Nós defendemos a justiça e a equidade.]
  \end{Phonetics}
\end{Entry}

\begin{Entry}{崇拜}{11,9}{⼭,⼿}
  \begin{Phonetics}{崇拜}{chong2bai4}[][HSK 6]
    \definition{v.}{adorar; idolatrar; venerar}
  \end{Phonetics}
\end{Entry}

\begin{Entry}{崇高}{11,10}{⼭,⾼}
  \begin{Phonetics}{崇高}{chong2gao1}[][HSK 7-9]
    \definition{adj.}{alto; elevado; sublime; nobre}
  \end{Phonetics}
\end{Entry}

%%%%%%%%%% 崖 %%%%%%%%%%
\subsection*{崖}\addcontentsline{loh}{figure}{崖}

\begin{Entry}{崖}{11}{⼭}
  \begin{Phonetics}{崖}{ya2}
    \definition{s.}{precipício | penhasco}
  \end{Phonetics}
\end{Entry}

%%%%%%%%%% 崛 %%%%%%%%%%
\subsection*{崛}\addcontentsline{loh}{figure}{崛}

\begin{Entry}{崛}{11}{⼭}
  \begin{Phonetics}{崛}{jue2}
    \definition{v.}{Literário: subir abruptamente; levantar"-se abruptamente; empinar}
  \end{Phonetics}
\end{Entry}

\begin{Entry}{崛起}{11,10}{⼭,⾛}
  \begin{Phonetics}{崛起}{jue2qi3}[][HSK 7-9]
    \definition{v.}{surgir abruptamente; subir repentinamente; aparecer de repente no horizonte | ganhar destaque; ascender}
  \end{Phonetics}
\end{Entry}

%%%%%%%%%% 崩 %%%%%%%%%%
\subsection*{崩}\addcontentsline{loh}{figure}{崩}

\begin{Entry}{崩}{11}{⼭}
  \begin{Phonetics}{崩}{beng1}
    \definition{v.}{colapsar |  estourar; quebrar | atingir por explosão | matar atirando; atirar; executar | (de um imperador) morrer | rachar; romper | atingir | executar atirando}
  \end{Phonetics}
\end{Entry}

\begin{Entry}{崩溃}{11,12}{⼭,⽔}
  \begin{Phonetics}{崩溃}{beng1kui4}[][HSK 7-9]
    \definition{v.}{colapsar; desmoronar; cair aos pedaços; as coisas estão destruídas; as emoções das pessoas estão fora de controle}
  \synonymref{倒闭}{dao3bi4}
  \synonymref{解体}{jie3ti3}
  \synonymref{破产}{po4/chan3}
  \antonymref{强大}{qiang2da4}
  \end{Phonetics}
\end{Entry}

%%%%%%%%%% 崭 %%%%%%%%%%
\subsection*{崭}\addcontentsline{loh}{figure}{崭}

\begin{Entry}{崭}{11}{⼭}
  \begin{Phonetics}{崭}{zhan3}
    \definition{adj.}{Literário: imponente; altaneiro; superior | Dialeto: bom; fresco; novo | excelente}
  \end{Phonetics}
\end{Entry}

\begin{Entry}{崭新}{11,13}{⼭,⽄}
  \begin{Phonetics}{崭新}{zhan3xin1}[][HSK 7-9]
    \definition{adj.}{novo em folha; completamente novo; muito novo}
  \synonymref{全新}{quan2xin1}
  \antonymref{报废}{bao4/fei4}
  \antonymref{陈旧}{chen2jiu4}
  \antonymref{破旧}{po4jiu4}
  \end{Phonetics}
\end{Entry}

%%%%%%%%%% 嵩 %%%%%%%%%%
\subsection*{嵩}\addcontentsline{loh}{figure}{嵩}

\begin{Entry}{嵩}{13}{⼭}
  \begin{Phonetics}{嵩}{song1}
    \definition*{s.}{Sobrenome: Song}
    \definition{adj.}{Arcaico: (montanhas) alto; elevado}
  \end{Phonetics}
\end{Entry}

\begin{Entry}{嵩山}{13,3}{⼭,⼭}
  \begin{Phonetics}{嵩山}{song1shan1}
    \definition*{s.}{Monte Song em Henan, montanha central das Cinco Montanhas Sagradas (五岳)}
  \seealsoref{五岳}{wu3yue4}
  \end{Phonetics}
\end{Entry}

%%%%%%%%%% 巅 %%%%%%%%%%
\subsection*{巅}\addcontentsline{loh}{figure}{巅}

\begin{Entry}{巅}{19}{⼭}
  \begin{Phonetics}{巅}{dian1}
    \definition[个]{s.}{pico da montanha; cume; topo da montanha}
  \end{Phonetics}
\end{Entry}

\begin{Entry}{巅峰}{19,10}{⼭,⼭}
  \begin{Phonetics}{巅峰}{dian1feng1}[][HSK 7-9]
    \definition{s.}{um cume; um pico de montanha}
  \end{Phonetics}
\end{Entry}

%%%%% EOF %%%%%

